
\section{Related work}
\label{sec:rh_related}
\begin{itemize}
\item \textbf{Realized variance and the HAR model.} Realized variance, the sum of squared high-frequency returns,
  makes volatility observable and forecastable \citep{AndersenBollerslevDieboldLabys2003,
  BarndorffNielsenShephard2002}, and the heterogeneous autoregressive (HAR) model of \citet{Corsi2009}, which
  regresses next-day realized variance on its averages over the past day, week, and month, has become the benchmark
  for forecasting it. Its extensions add jumps \citep{AndersenBollerslevDiebold2007}, signed semivariances
  \citep{PattonSheppard2015}, leverage \citep{CorsiReno2012}, and the measurement error of realized variance
  \citep{BollerslevPattonQuaedvlieg2016}, and the panel HAR of \citet{BollerslevHoodHussLebedev2018} estimates a
  single pooled regression across a cross-section of assets. These are models of daily variance. We carry the HAR
  structure to the intraday interval, pool it across the intervals of one index rather than across assets, and ask
  how much a large set of exogenous predictors adds to it.
\item \textbf{Pooled intraday forecasting and overnight information.} Intraday volatility follows a pronounced
  U-shaped pattern across the trading session, and scaling intraday returns or realized variances by a diurnal
  factor is standard practice \citep{AndersenBollerslev1997,TaylorXu1997}; \citet{AndersenSuTodorovZhang2024}
  estimate this curve for the S\&P~500 futures around the clock and find that volatility outside U.S. trading hours
  matters more in the later part of their sample. With this scaling, one model can be estimated across the intervals
  of the regular session: \citet{EngleSokalska2012} estimate a multiplicative component GARCH model of intraday
  volatility, with daily, diurnal, and stochastic intraday components, pooled across a large cross-section of
  stocks, and \citet{ZhangZhangCucuringuQian2024} fit machine-learning models of realized volatility across all
  intraday buckets of the day, with predictors built from realized volatility alone. Both leave the overnight return
  out of estimation, and so do we (Section~\ref{sec:rh_methods}). Yet information from outside the session helps
  forecast it: after-hours realized variance \citep{ChenYuZivot2012} and overnight information, including the
  realized volatility of related assets that trade while the home market is closed \citep{TodorovaSoucek2014},
  improve volatility forecasts, and for the S\&P~500 futures \citet{StroudJohannes2014} argue that deleting the
  overnight returns removes important information. These studies forecast daily or session volatility; we carry the
  idea to the intraday interval by computing our predictors over the futures' full intraday grid, overnight
  intervals included (Section~\ref{sec:rh_predictors}).
\item \textbf{Machine learning for realized volatility.} Applications of bagging, the lasso, boosting, random
  forests, and neural networks to realized volatility \citep{HillebrandMedeiros2010,AudrinoKnaus2016,
  MittnikRobinzonovSpindler2015,LuongDokuchaev2018,RahimikiaPoon2020} preceded the comparison of these methods with
  the HAR model for daily realized volatility in \citet{christensen2023machine}, much as \citet{gu2020empirical}
  compared linear, penalized, tree-based, and neural-network models for returns; \citet{BrancoRubesamZevallos2024}
  find that nonlinear models rarely beat linear ones. Most of this work targets daily volatility; we compare
  penalized linear models with tree ensembles on every regular-hours interval.
\item \textbf{Dense or sparse signals.} In macroeconomic, microeconomic, and financial data, the posterior over
  predictive models does not typically concentrate on a single sparse model \citep{GiannoneLenzaPrimiceri2021}, and
  a characteristics-sparse stochastic discount factor cannot adequately summarize the cross-section of expected
  stock returns \citep{KozakNagelSantosh2020}; \citet{ShenXiu2025} ask whether machine-learning methods can learn
  weak signals. That our exogenous signal is dense and weak, spread thinly over many series, is a hypothesis we
  test, not a premise: ridge regression suits a dense signal and the lasso a sparse one, and we compare penalized
  linear models with tree ensembles across the number of exogenous series and the length of the training window
  (Section~\ref{sec:rh_fi}). Unlike the complex models of \citet{KellyMalamudZhou2024}, which have more parameters
  than observations, ours have far fewer predictors than training observations.
\item \textbf{Alternative data.} Sentiment, attention, and news measures have mostly been used to forecast daily
  realized volatility \citep{AudrinoSigristBallinari2020,CaporinPoli2017}; at the intraday horizon,
  \citet{BehrendtSchmidt2018} relate Twitter sentiment and activity to the intraday volatility of individual stocks.
  We ask what such data add to a pooled intraday model of the index's realized variance, next to the futures' own
  realized moments, the trading of the index's constituents, implied volatility, and the FOMC calendar.
\item \textbf{Economic evaluation of volatility forecasts and 0DTE options.} A long tradition ranks volatility
  forecasts by the profits of trading options on them, buying (selling) straddles when the price implied by the
  forecast exceeds (falls short of) the market price \citep{EngleHongKaneNoh1993,NohEngleKane1994,
  BandiRussellYang2008}; statistical gains need not survive transaction costs
  \citep{HarveyWhaley1992,ChanJhaKalimipalli2009}. Because the variance risk premium is positive on average
  \citep{CarrWu2009,BollerslevTauchenZhou2009}, zero-beta at-the-money straddles and delta-hedged index options lose
  money on average \citep{CovalShumway2001,BakshiKapadia2003}, which makes the unconditional short straddle the
  natural benchmark for a strategy that trades straddles on a variance forecast. A recent literature studies the
  pricing, trading, and volatility effects of zero-days-to-expiration (0DTE) options on the S\&P~500
  \citep{BandiFusariReno2023,BeckmeyerBrangerGayda2023,DimErakerVilkov2024}, and \citet{ChongTodorov2026} ask
  whether equity and 0DTE options markets agree about volatility and trade their disagreement. Nearest to our market
  and bar length, \citet{MaischZhangUlrichZimmer2026} predict the half-hourly delta-hedged returns of S\&P~500 index
  options from contract characteristics, Greeks, and option-implied measures of the index's return distribution,
  and sort the contracts into long--short portfolios. Our forecasts belong to the realized volatility literature,
  and options enter only as an application: one 0DTE at-the-money straddle a day through the closing half hour,
  long or short as the forecast for that interval exceeds or falls short of the implied variance
  (Section~\ref{sec:rh_econ}).
\end{itemize}

